\documentclass[10pt]{article}

\usepackage[utf8]{inputenc}

\usepackage[T1]{fontenc}

\usepackage{amsmath}

\usepackage{amsfonts}

\usepackage{amssymb}

\usepackage[version=4]{mhchem}

\usepackage{stmaryrd}

\usepackage{bbold}

\usepackage{mathrsfs}



\begin{document}

где $\varphi(x)$ - приближающая функция, $R(x)$ - погрешность, относительно к-рой в общем случае известно только, что



\[

\lim _{x \rightarrow \infty}(R(x) / \varphi(x))=0 .

\]



Краткая запись: $f(x)=\varphi(x)+o \quad(\varphi(x))$ пли $f(x) \sim \varphi(x)$ (см. Асимптотическая формула).



Нахождение Ч. ф. а. является одной из важнейших задач аналитич. теории чисел. Это объясняется тем, что почти все числовые функции с интересными арифметич. свойствами характеризуются крайней неправильностью своего изменения при возрастании аргумента. При рассмотрении вместо числовоїі функции $f(x)$ ее среднего значения $\frac{1}{x} \sum_{n \leqslant x} f(n)(n-$ натуральное) «неправильность» функции $f(x)$ сглаживается. Поэтому типичной для числовых функций является задача получения асимптотики для их средних значений. Напр., среднее значение функции, означающей число делителей $n$, будет равняться



\[

\frac{1}{n} \sum_{m \leqslant n} \tau(n) \sim \ln n .

\]



Возникающая здесь проблема наилучшей возможної оценки погренности в асимптотич. равенстве для многих функций, в частности для функций $\tau(n)$, в 1984 еще не решена (см. Аналитическая теория чисел).



Не менее важную роль играет Ч. ф. а. при решении аддитивных проблем. Для многих из них неизвестны прямые доказательства существования разложения целого числа $n$ на слагаемые заданного вида. Однако, как только удается вывести асимптотику для числа искомых разложений, отсюда уже следует существование требуемого разложения для ліобого достаточно большого числа $n$.



ЧИСТЫЙ ПОДМОДУЛЬ в с м ы с л $\mathrm{e}$ К о н а такой подмодуль $A$ правого $R$-модуля $B$, что для любого левого $R$-модуля $C$ естественный гомоморфизм абелевых групп



\[

A \otimes_{R} C \rightarrow B \otimes_{R} C

\]



инъективен. Это эквивалентно следующему условию: если система уравнений



\[

\sum_{i=1}^{n} x_{i} \lambda_{i j}=a_{j} \text {, где } 1 \leqslant j \leqslant m, \lambda_{i j} \in R, a_{j} \in A,

\]



имеет решение в $B$, то она имеет решение и в $A$ (ср. Плоский модуль). Ч. п. является любое прямое слагаемое. Любой подмопуль любого правого $R$-модуля чист тогда и только тогда, когда $R$ - регулярное кольцо .



В случае абелевых групп (т. е. при $R=\mathbb{Z}$ ) әквивалентны следующие утверждения: (1) $A-$ ч и с т а я (или с е р в а н т н а я) подгруппа в $B$; (2) $n A=A \cap n B$ для любого натурального $n$; (3) $A / n A-$ прямое слагаемое в $B / n A$ для любого натурального $n ;(4)$ если $C \subseteq A$ и $A / C$ - конечно порожденная группа, то $A / C-$ прямое слагаемое в $B / C ;(5)$ каждый смежный класс факторгрупшы $B / A$ содержит әлемент того же порядка, что и порядок этого смежного класса; (6) если $A \subseteq C \subseteq B$ и $C / A$ конечно порождена, то $A-$ прямое слагаемое в $C$. Если выполнение свойства (2) требуется лишь для простых $n$, то $A$ наз. с л а б о с е р в а н т н о ї по дг р у П п о й.



Аксиоматич. подход к понятию чистоты базируется на рассмотрении класса мономорфизмов Ћ' ного следующим требованиям (здесь $A \subseteq \omega B$ означает, что $A$ - подмодуль в $B$ и естественное вложение принадлежит $\left.\mathscr{\Re}_{\omega}\right):\left(Ч 0^{\prime}\right)$ если $A$-прямое слагаемое в $B$, то $A \subseteq \omega B$; (Ч1 $\left.{ }^{\prime}\right)$ если $A \subseteq{ }_{\omega} B$ и $B \subseteq{ }_{\omega} C$, то $A \subseteq{ }_{\omega} C$; (ч2') если $A \subseteq B \subseteq C$ и $A \subseteq \omega C$, то $A \subseteq \omega B$; (ЧЗ $3^{\prime}$ ) если $A \sqsubseteq{ }_{\omega} B$ и $K \sqsubseteq A$, то $A / K \sqsubseteq{ }_{\omega} B / K$; (ㄴ') ес ли $K \sqsubseteq A \subseteq B$, $K \sqsubseteq{ }_{\omega} B$ и $A / K \sqsubseteq{ }_{\omega} B / K$, то $A \subseteq а . B$. Рассмотрение класса Л̆ к относительной гомологич. алгебре. Напр., модуль $Q$ наз. $\omega$-и н ъ е к т и в н м м, если из $A \subseteq \omega B$ вытекает, что всякий гомоморфизм из $A$ в $Q$ может быть продолжен до гомоморфизма из $B$ в $Q$ (cp. Инъективный модуль). Сервантно инъективная абелева группа наз. а л г е б ра и ч е с к и ко м п а к тной. Эквивалентны следующие свойства абелевой группы $Q:$ (1) $Q$ алгебраически компактна; (2) $Q$ выделяется прямым слагаемым из любой группы, содержащей ее в качестве сервантной подгруппы; (3) $Q$ является прямым слагаемым группы, допускающей компактную топологию; (4) система уравнений над $Q$ разрешима, если разрепима каждая пз ее конечных подсистем.



Лum.: [1] М и ши и а А. П., С кор н я ко в Л. А., Абелевы группы и модули, М., $1969 ;$ [2] С к л я р е н к о Е. Г, "Успехи матем. наук", 1978, т. 33 , в. 3 , с. 85-120;

Ф е й с К., Алгебра: кольца, модули и категории, пер. с англ., Т. $1-2$, М., $1977-79 ;[4]$, $Ф$ у к с Ј., Бесконечные абелевы грушпы, пер. с англ., т. 1-2, М., $1974-77$.



Л. А. Скорняков.





\end{document}